% GENERATED FILE. Edit canonical lessons, then run npm run generate:books.
% canonical-source-hash: fnv1a64:3ebe730d8ec4f482
% canonical-lessons: FR-C476-inquietude, FR-C476-anxieux, FR-C476-mefier, FR-C476-optimiste, FR-C476-pessimiste

\chapter{Optimistes et inquiets}
\label{ch:fr-optimistes-et-inquiets}
\hlchaptermodality{476}

\begin{chapteropening}
\textbf{By the end of this chapter:} \emph{I can say whether someone worries or looks on the bright side.}

Complete the last lesson of chapter 476: I can say whether someone worries or looks on the bright side.
\end{chapteropening}

\section[l'inquiétude]{l'inquiétude — worry, anxiety}

\label{lesson:FR-C476-inquietude}



\begin{quote}
Before the new one: say the French for resentful, then the French for to reconcile.
\end{quote}

\subsection*{You'll want to know: l'inquiétude}
\textbf{l'inquiétude} — \textquotedblleft{}worry, anxiety\textquotedblright{}.

\begin{quote}
Elle attend le résultat avec inquiétude.
\end{quote}

\begin{quote}
\emph{She is waiting for the result with worry.}
\end{quote}

Built on \emph{inquiet} \textquotedblleft{}worried\textquotedblright{}. Feminine: \emph{une inquiétude}.

\begin{cousinweb}
\emph{l'inquiétude} is a learned borrowing from Latin \emph{inquiētūdō} \textquotedblleft{}restlessness\textquotedblright{} (EtymDB), built on \emph{quiēs} \textquotedblleft{}rest\textquotedblright{} with the negative \emph{in-} (Lewis \& Short).
\end{cousinweb}

\begin{grammarlens}[title={What you've built}]
One more word for telling, explaining and describing at B1.
\end{grammarlens}

\subsection*{Your turn}
\begin{itemize}
\raggedright
  \item \emph{Say it:} \emph{l'inquiétude}
  \item \emph{Say it:} \emph{l'inquiétude}, once more
  \item \emph{Say it:} a sentence of your own with \emph{l'inquiétude}
  \item \emph{Recall:} say the French for resentful, then the French for to reconcile, then say \emph{l'inquiétude} again
\end{itemize}

\begin{tcolorbox}[breakable,colback=teal!4,colframe=teal!35!black,title={Before you move on}]
What does l'inquiétude mean? (\textquotedblleft{}Worry, anxiety\textquotedblright{}.) Say it once more.
\end{tcolorbox}
\section[anxieux]{anxieux — anxious}

\label{lesson:FR-C476-anxieux}



\begin{quote}
Before the new one: say the French for to reconcile, then the French for worry.
\end{quote}

\subsection*{You'll want to know: anxieux}
\textbf{anxieux} — \textquotedblleft{}anxious\textquotedblright{}.

\begin{quote}
Avant un examen, je suis toujours anxieux.
\end{quote}

\begin{quote}
\emph{Before an exam I am always anxious.}
\end{quote}

Feminine \emph{anxieuse}. Stronger than \emph{inquiet}.

\begin{cousinweb}
\emph{anxieux} is a learned borrowing from Latin \emph{anxius} \textquotedblleft{}troubled\textquotedblright{} (EtymDB), from \emph{angere} \textquotedblleft{}to squeeze, to distress\textquotedblright{} (Lewis \& Short).
\end{cousinweb}

\begin{grammarlens}[title={What you've built}]
One more word for telling, explaining and describing at B1.
\end{grammarlens}

\subsection*{Your turn}
\begin{itemize}
\raggedright
  \item \emph{Say it:} \emph{anxieux}
  \item \emph{Say it:} \emph{anxieux}, once more
  \item \emph{Say it:} a sentence of your own with \emph{anxieux}
  \item \emph{Recall:} say the French for to reconcile, then the French for worry, then say \emph{anxieux} again
\end{itemize}

\begin{tcolorbox}[breakable,colback=teal!4,colframe=teal!35!black,title={Before you move on}]
What does anxieux mean? (\textquotedblleft{}Anxious\textquotedblright{}.) Say it once more.
\end{tcolorbox}
\section[méfier]{méfier — to be wary, to distrust (se méfier)}

\label{lesson:FR-C476-mefier}



\begin{quote}
Before the new one: say the French for worry, then the French for anxious.
\end{quote}

\subsection*{You'll want to know: méfier}
\textbf{méfier} — \textquotedblleft{}to be wary, to distrust (se méfier)\textquotedblright{}.

\begin{quote}
Je me méfie des gens trop gentils.
\end{quote}

\begin{quote}
\emph{I am wary of people who are too nice.}
\end{quote}

Always reflexive in use: \emph{se méfier de}. The adjective is \emph{méfiant}.

\begin{cousinweb}
\emph{se méfier} is built from \emph{mé-} \textquotedblleft{}badly\textquotedblright{} and \emph{se fier} \textquotedblleft{}to trust\textquotedblright{}, which goes back to Vulgar Latin \emph{fīdāre}, from Latin \emph{fīdere} \textquotedblleft{}to trust\textquotedblright{} (Lewis \& Short).
\end{cousinweb}

\begin{grammarlens}[title={What you've built}]
One more word for telling, explaining and describing at B1.
\end{grammarlens}

\subsection*{Your turn}
\begin{itemize}
\raggedright
  \item \emph{Say it:} \emph{méfier}
  \item \emph{Say it:} \emph{méfier}, once more
  \item \emph{Say it:} a sentence of your own with \emph{méfier}
  \item \emph{Recall:} say the French for worry, then the French for anxious, then say \emph{méfier} again
\end{itemize}

\begin{tcolorbox}[breakable,colback=teal!4,colframe=teal!35!black,title={Before you move on}]
What does méfier mean? (\textquotedblleft{}To be wary, to distrust (se méfier)\textquotedblright{}.) Say it once more.
\end{tcolorbox}
\section[optimiste]{optimiste — optimistic; optimist}

\label{lesson:FR-C476-optimiste}



\begin{quote}
Before the new one: say the French for anxious, then the French for to be wary.
\end{quote}

\subsection*{You'll want to know: optimiste}
\textbf{optimiste} — \textquotedblleft{}optimistic; optimist\textquotedblright{}.

\begin{quote}
Malgré la pluie, elle reste optimiste.
\end{quote}

\begin{quote}
\emph{In spite of the rain, she stays optimistic.}
\end{quote}

The same form for a man and a woman.

\begin{cousinweb}
\emph{optimiste} is built on \emph{optimisme}, a modern learned word made from Latin \emph{optimus} \textquotedblleft{}best\textquotedblright{} (EtymDB; Lewis \& Short).
\end{cousinweb}

\begin{grammarlens}[title={What you've built}]
One more word for telling, explaining and describing at B1.
\end{grammarlens}

\subsection*{Your turn}
\begin{itemize}
\raggedright
  \item \emph{Say it:} \emph{optimiste}
  \item \emph{Say it:} \emph{optimiste}, once more
  \item \emph{Say it:} a sentence of your own with \emph{optimiste}
  \item \emph{Recall:} say the French for anxious, then the French for to be wary, then say \emph{optimiste} again
\end{itemize}

\begin{tcolorbox}[breakable,colback=teal!4,colframe=teal!35!black,title={Before you move on}]
What does optimiste mean? (\textquotedblleft{}Optimistic; optimist\textquotedblright{}.) Say it once more.
\end{tcolorbox}
\section[pessimiste]{pessimiste — pessimistic; pessimist}

\label{lesson:FR-C476-pessimiste}



\begin{quote}
Before the new one: say the French for to be wary, then the French for optimistic.
\end{quote}

\subsection*{You'll want to know: pessimiste}
\textbf{pessimiste} — \textquotedblleft{}pessimistic; pessimist\textquotedblright{}.

\begin{quote}
Mon oncle est pessimiste : il pense que tout va mal finir.
\end{quote}

\begin{quote}
\emph{My uncle is pessimistic: he thinks everything will end badly.}
\end{quote}

The same form for a man and a woman.

\begin{cousinweb}
\emph{pessimiste} is built on \emph{pessimisme}, a modern learned word made from Latin \emph{pessimus} \textquotedblleft{}worst\textquotedblright{} (EtymDB; Lewis \& Short).
\end{cousinweb}

\begin{grammarlens}[title={What you've built}]
One more word for telling, explaining and describing at B1.
\end{grammarlens}

\subsection*{Your turn}
\begin{itemize}
\raggedright
  \item \emph{Say it:} \emph{pessimiste}
  \item \emph{Say it:} \emph{pessimiste}, once more
  \item \emph{Say it:} a sentence of your own with \emph{pessimiste}
  \item \emph{Recall:} say the French for to be wary, then the French for optimistic, then say \emph{pessimiste} again
\end{itemize}

\begin{tcolorbox}[breakable,colback=teal!4,colframe=teal!35!black,title={Before you move on}]
What does pessimiste mean? (\textquotedblleft{}Pessimistic; pessimist\textquotedblright{}.) Say it once more.
\end{tcolorbox}
